% !TEX root = ../main.tex
\section{Ejercicio propuesto 2: comparación con heapq y aplicación en Dijkstra}

\subsection{Protocolo experimental y configuración del entorno}
Los experimentos fueron ejecutados bajo Python 3.12 sobre un entorno controlado de 64 bits con la semilla fija $\texttt{SEMILLA\_GLOBAL} = 2026$. Cada medición de tiempo utilizó la función de alta precisión \texttt{time.perf\_counter\_ns} aislando la inicialización de estructuras fuera de los bucles cronometrados para eliminar perturbaciones de entrada/salida.

% -------------------------------------------------------------------------
\subsection{Escenario A: operaciones básicas de inserción y extracción}
Se evaluaron secuencias de $n \in \{1\,000, 5\,000, 10\,000, 25\,000\}$ elementos generados pseudoaleatoriamente. Se realizaron 10 repeticiones independientes por configuración, comprobando en cada una de ellas que ambas estructuras devuelven exactamente la misma secuencia ordenada de claves.

\begin{table}[H]
    \caption{Comparativa del Escenario A (Inserción y Extracción sobre 10 repeticiones)}
    \label{tab:escenario_a}
    \centering
    \scriptsize
    \begin{tabular*}{\textwidth}{@{\extracolsep{\fill}}rcccccccc}
        \toprule
        & \multicolumn{4}{c}{\textbf{Inserción de $n$ Claves (ms)}} & \multicolumn{4}{c}{\textbf{Extracción Completa de $n$ Mínimos (ms)}} \\
        \cmidrule(lr){2-5} \cmidrule(lr){6-9}
        $\mathbf{n}$ & \textbf{Heapq (Med)} & \textbf{Heapq (Min-Max)} & \textbf{Fib (Med)} & \textbf{Fib (Min-Max)} & \textbf{Heapq (Med)} & \textbf{Heapq (Min-Max)} & \textbf{Fib (Med)} & \textbf{Fib (Min-Max)} \\
        \midrule
        1\,000  & 0.59 & 0.54 -- 1.99  & 0.86  & 0.77 -- 1.40  & 1.17  & 1.09 -- 1.83  & 13.55  & 12.79 -- 26.57 \\
        5\,000  & 4.91 & 3.67 -- 7.57  & 4.60  & 4.05 -- 29.87 & 11.00 & 8.45 -- 19.96 & 98.58  & 91.37 -- 110.63 \\
        10\,000 & 7.48 & 7.06 -- 27.69 & 9.71  & 9.11 -- 41.32 & 20.16 & 17.49 -- 40.15 & 224.94 & 206.23 -- 335.13 \\
        25\,000 & 22.89 & 19.66 -- 61.20 & 29.57 & 23.14 -- 69.05 & 63.30 & 56.31 -- 85.70 & 654.74 & 627.52 -- 840.95 \\
        \bottomrule
    \end{tabular*}
\end{table}

\begin{figure}[H]
    \centering
    \includegraphics[width=0.95\textwidth]{fig_escenario_a_tiempos.png}
    \caption{Tiempos de ejecución en el Escenario A (Inserción y Extracción frente a $n$)}
    \label{fig:escenario_a}
\end{figure}

\subsubsection*{Hallazgo experimental}
En inserción, los tiempos de ambas estructuras son comparables (en $n=5\,000$, Fibonacci fue incluso más veloz con 4.60 ms frente a 4.91 ms debido a su inserción perezosa en tiempo estrictamente constante sin reordenar). Sin embargo, en extracción completa, \texttt{heapq} resulta significativamente superior debido a su implementación en lenguaje C sobre contigüidad de memoria en arreglos, mientras que el montículo de Fibonacci incurre en la sobrecarga de reconstruir enlaces y punteros de objetos dinámicos durante la consolidación.

% -------------------------------------------------------------------------
\subsection{Escenario B: carga intensiva de disminuciones de clave}
Para contrastar el costo asintótico de \texttt{decrease-key} ($O(1)$ en Fibonacci frente a $O(\log n)$ en montículos binarios), se insertaron $n$ elementos y se ejecutaron $q = 10n$ operaciones válidas de disminución de clave, intercalando una extracción cada 20 disminuciones.

\begin{table}[H]
    \caption{Resultados del Escenario B: $q = 10n$ disminuciones con extracciones intercaladas}
    \label{tab:escenario_b}
    \centering
    \scriptsize
    \begin{tabular*}{\textwidth}{@{\extracolsep{\fill}}rccccccccccc}
        \toprule
        $\mathbf{n}$ & $\mathbf{q}$ & $\mathbf{T_{\text{heapq}}}$ (ms) & $\mathbf{T_{\text{fib}}}$ (ms) & \textbf{Físico} & \textbf{Activos} & \textbf{Cortes} & \textbf{Cascadas} & \textbf{Enlaces} & $\mathbf{t(H)}$ & $\mathbf{m(H)}$ & $\mathbf{\Phi(H)}$ \\
        \midrule
        1\,000  & 10\,000  & 16.23  & 15.79  & 5\,563   & 500   & 0 & 0 & 4\,696   & 6 & 0 & 6 \\
        5\,000  & 50\,000  & 118.70 & 107.51 & 27\,514  & 2\,500 & 2 & 0 & 29\,348  & 5 & 0 & 5 \\
        10\,000 & 100\,000 & 316.29 & 223.45 & 55\,322  & 5\,000 & 1 & 0 & 63\,963  & 5 & 0 & 5 \\
        20\,000 & 200\,000 & 767.94 & 568.47 & 110\,220 & 10\,000 & 7 & 0 & 138\,097 & 5 & 0 & 5 \\
        \bottomrule
    \end{tabular*}
\end{table}

\begin{figure}[H]
    \centering
    \begin{subfigure}[b]{0.485\textwidth}
        \centering
        \includegraphics[width=\textwidth]{fig_escenario_b_disminuciones.png}
        \caption{Tiempo de cómputo frente a $q$ disminuciones}
        \label{fig:esc_b_tiempo}
    \end{subfigure}
    \hfill
    \begin{subfigure}[b]{0.485\textwidth}
        \centering
        \includegraphics[width=\textwidth]{fig_escenario_b_tamano_heapq.png}
        \caption{Crecimiento de entradas obsoletas en \texttt{heapq}}
        \label{fig:esc_b_tamano}
    \end{subfigure}
    \caption{Evaluación del comportamiento en el Escenario B bajo alta tasa de disminución}
    \label{fig:escenario_b}
\end{figure}

\subsubsection*{Ventaja práctica del Montículo de Fibonacci}
Como se visualiza en la Figura \ref{fig:escenario_b}, el montículo de Fibonacci **supera de forma concluyente a \texttt{heapq}** cuando el volumen de disminuciones de clave es dominante ($q \ge 10n$):
\begin{itemize}
    \item Para $n = 20\,000$ y $q = 200\,000$, el montículo de Fibonacci resolvió la carga en **568.47 ms**, mientras que \texttt{heapq} requirió **767.94 ms** (una reducción de tiempo del $26\%$).
    \item La Figura \ref{fig:esc_b_tamano} revela el motivo físico: la técnica de invalidación perezosa provocó que la lista física interna de \texttt{heapq} acumulara **110\,220 entradas**, a pesar de existir únicamente **10\,000 elementos activos**. El 91\% del espacio correspondió a basura informática que degrada las operaciones posteriores.
\end{itemize}

% -------------------------------------------------------------------------
\subsection{Escenario C: algoritmo de Dijkstra en grafos sintéticos}
Se implementó el algoritmo de caminos mínimos de Dijkstra acoplado a ambas estructuras de prioridad. Se evaluaron grafos dispersos ($|E| \approx 3|V|$) y de densidad intermedia ($|E| \approx 10|V|$). En todas las instancias se comprobó la igualdad estricta de las distancias devueltas ($\texttt{dist\_heapq} = \texttt{dist\_fib}$).

\begin{table}[H]
    \caption{Tiempos de ejecución de Dijkstra sobre grafos de distinta densidad}
    \label{tab:dijkstra}
    \centering
    \small
    \begin{tabular*}{\textwidth}{@{\extracolsep{\fill}}lccccc}
        \toprule
        \textbf{Densidad} & $\mathbf{|V|}$ & $\mathbf{|E|}$ & $\mathbf{T_{\text{heapq}}}$ (ms) & $\mathbf{T_{\text{fib}}}$ (ms) & \textbf{Concordancia} \\
        \midrule
        Disperso ($3|V|$)    & 500   & 1\,500  & 3.18  & 8.02   & 100\% Idéntico \\
        Disperso ($3|V|$)    & 1\,000 & 3\,000  & 5.61  & 19.34  & 100\% Idéntico \\
        Disperso ($3|V|$)    & 2\,000 & 6\,000  & 21.36 & 43.60  & 100\% Idéntico \\
        Disperso ($3|V|$)    & 4\,000 & 12\,000 & 32.10 & 128.03 & 100\% Idéntico \\
        \midrule
        Intermedio ($10|V|$) & 500   & 5\,000  & 6.94  & 10.86  & 100\% Idéntico \\
        Intermedio ($10|V|$) & 1\,000 & 10\,000 & 30.10 & 38.40  & 100\% Idéntico \\
        Intermedio ($10|V|$) & 2\,000 & 20\,000 & 42.79 & 86.51  & 100\% Idéntico \\
        Intermedio ($10|V|$) & 4\,000 & 40\,000 & 75.22 & 120.07 & 100\% Idéntico \\
        \bottomrule
    \end{tabular*}
\end{table}

\begin{figure}[H]
    \centering
    \includegraphics[width=0.95\textwidth]{fig_escenario_c_dijkstra.png}
    \caption{Comparación de tiempos en Dijkstra según densidad del grafo}
    \label{fig:escenario_c}
\end{figure}

\subsubsection*{Validación de casos frontera en Dijkstra}
Se evaluaron exitosamente los siguientes casos atípicos:
\begin{itemize}
    \item \textbf{Grafos desconectados:} Se validó que para vértices inalcanzables ambas implementaciones devuelven rigurosamente $\text{dist}[u] = \infty$.
    \item \textbf{Aristas de peso cero:} Se procesaron aristas de costo nulo confirmando la preservación de caminos mínimos sin bucles infinitos.
    \item \textbf{Pesos negativos:} El código rechaza formalmente aristas con pesos negativos lanzando la excepción \texttt{ValueError}, dado que Dijkstra presupone costos no negativos.
\end{itemize}

% -------------------------------------------------------------------------
\subsection{Respuestas a las preguntas de análisis de la guía}

\subsubsection*{1. ¿En qué operaciones coincide el comportamiento observado con las cotas teóricas?}
El comportamiento coincide con las cotas asintóticas en las operaciones de inserción ($O(1)$ constante) y unión ($O(1)$), las cuales no dependen de la escala del montículo. En extracción, la cota $O(\log n)$ se manifiesta en el crecimiento logarítmico del tiempo, mientras que en \texttt{decrease-key} el costo amortizado $O(1)$ permitió que para 200\,000 operaciones el tiempo por disminución fuera de tan solo $2.84\,\mu\text{s}$.

\subsubsection*{2. ¿Por qué heapq puede ser más rápido aunque decrease-key sea $O(\log n)$?}
Porque el módulo \texttt{heapq} de la biblioteca estándar de Python está compilado internamente en **lenguaje C optimizado**. Opera sobre contigüidad física de memoria en listas nativas, lo que maximiza la tasa de aciertos en la memoria caché del procesador (\textit{cache locality}) y elimina la creación de objetos dinámicos de Python y la manipulación de punteros circulares que ralentizan al intérprete.

\subsubsection*{3. ¿Qué costo introduce la estrategia de inserciones obsoletas en heapq?}
Introduce un doble sobrecosto:
\begin{itemize}
    \item \textbf{Espacial:} La memoria física requerida se dispara proporcionalmente a la cantidad de disminuciones. En el Escenario B, el arreglo acumuló 110\,220 entradas para albergar únicamente 10\,000 claves activas ($11\times$ de inflación de memoria).
    \item \textbf{Temporal:} El tamaño agrandado del arreglo degrada las operaciones posteriores de \texttt{heappop} y \texttt{heappush}, puesto que la altura del árbol binario pasa de $\log_2(n)$ a $\log_2(n + q)$.
\end{itemize}

\subsubsection*{4. ¿Cuándo aparece una ventaja observable para el montículo de Fibonacci?}
Aparece cuando la carga computacional está fuertemente desbalanceada hacia las disminuciones de clave frente a las extracciones (relación $q/n \ge 10$). En tales circunstancias, el costo amortizado $O(1)$ de Fibonacci ahorra el sobrecosto logarítmico repetido de \texttt{heapq} y evita la inflación del tamaño físico del montículo, logrando superar a la biblioteca en C (568 ms vs. 767 ms en nuestras pruebas).

\subsubsection*{5. ¿Cómo cambia $\Phi(H)$ durante inserciones, consolidaciones y cortes?}
\begin{itemize}
    \item \textbf{Inserción:} Agrega un árbol a la lista de raíces, por lo que $\Delta \Phi = +1$.
    \item \textbf{Consolidación:} Si se reducen $k$ raíces mediante enlaces, $\Delta t = -(k - 1)$, disminuyendo fuertemente el potencial ($\Delta \Phi < 0$) para compensar el costo del recorrido.
    \item \textbf{Corte simple:} Añade una raíz y marca un nodo padre, por lo que $\Delta \Phi = +1 + 2 = +3$.
    \item \textbf{Corte en cascada:} Cada nodo en cascada pasa a ser raíz ($+1$) y pierde su marca ($-2$), produciendo un cambio neto negativo $\Delta \Phi = 1 - 2 = -1$.
\end{itemize}

\subsubsection*{6. ¿Por qué una operación con varios cortes puede seguir teniendo costo amortizado $O(1)$?}
Sea una operación de \texttt{decrease-key} que desencadena $c$ cortes en cascada. El costo real es $c_i = O(c)$. El primer corte genera $\Delta \Phi \le 3$, pero cada uno de los $c - 1$ cortes en cascada subsecuentes desmarca un nodo y crea una raíz, aportando $\Delta \Phi = -1$. El cambio total de potencial es:
\begin{equation}
    \Delta \Phi \le 3 - (c - 1) = 4 - c
\end{equation}
Sumando el costo real y la variación de potencial:
\begin{equation}
    \hat{c}_i = c_i + \Delta \Phi = O(c) - c + 4 = O(1)
\end{equation}
El potencial previamente almacenado en los nodos marcados absorbe exactamente el costo lineal de la cascada, preservando la cota constante amortizada.

\subsubsection*{7. ¿Qué errores de implementación detectaron los verificadores de invariantes?}
El método \texttt{validate()} permitió detectar durante el desarrollo:
\begin{itemize}
    \item Nodos promovidos a raíz que conservaban erróneamente su referencia a \texttt{parent} o mantenían la bandera $\text{mark} = \text{True}$.
    \item Desajustes en el grado \texttt{degree} de nodos padres al desvincular hijos en cortes simples.
    \item Asignaciones incorrectas de punteros en la desconexión de listas circulares de un solo elemento ($x.\text{right} == x$), que producían ciclos infinitos.
\end{itemize}

\subsubsection*{8. ¿Recomendaría esta implementación para un sistema real en Python? Sustente la decisión}
**No se recomienda** para sistemas de producción estándar en Python interpretado. Aunque el montículo de Fibonacci posee una cota asintótica superior para grafos muy densos, la sobrecarga por manejo de objetos dinámicos, referencias de punteros y recolección de basura en CPython introduce constantes multiplicativas tan altas que el módulo \texttt{heapq} en lenguaje C resulta más eficiente y confiable en la gran mayoría de las aplicaciones prácticas de la industria.
